% ============================================================================
%  CAPÍTULO III: DESARROLLO DEL TRABAJO (parte 2)
% ============================================================================

\section{Ejecución de la migración}

\subsection{Conteos previos y respaldo del origen}

La ejecución comenzó con la verificación de los servicios y el conteo de las
8 tablas del origen. El script \texttt{02\_migrate\_and\_backup.sh} esperó a
que \texttt{db-origen} y \texttt{db-destino} estuvieran \emph{healthy} y a que
el migrador estuviera en ejecución. Luego generó un respaldo del origen con
\texttt{mysqldump --single-transaction --routines --triggers --events},
comprimido con \texttt{gzip -9}. El archivo resultante fue
\texttt{backups/origin\_20261120\_080940.sql.gz}, de 32 KB, con SHA-256
\texttt{6a2f3a958c4732743302299eb7e4f5d722924f749dbf6fa9764a2936428be3b3}.
La Tabla~\ref{tab:conteos-previos} reproduce los conteos.

\begin{table}[H]
  \centering
  \caption{Conteos del origen antes de la migración}
  \label{tab:conteos-previos}
  \contenidoConFuente{%
    \begin{tablaAPA}
      \begin{tabular}{@{}L{6.4cm}R{3.0cm}@{}}
        \toprule
        \textbf{Tabla} & \textbf{Filas} \\
        \midrule
        \texttt{sucursales} & 5 \\
        \texttt{socios} & 55 \\
        \texttt{cuentas\_ahorro} & 66 \\
        \texttt{tipos\_credito} & 5 \\
        \texttt{solicitudes\_credito} & 50 \\
        \texttt{prestamos} & 32 \\
        \texttt{cuotas\_amortizacion} & 516 \\
        \texttt{transacciones} & 740 \\
        \midrule
        \textbf{Total} & \textbf{1\,469} \\
        \bottomrule
      \end{tabular}
    \end{tablaAPA}%
  }{Elaboración propia con datos de \texttt{evidence/02\_ejecucion\_migracion.txt}.}
\end{table}

\subsection{Ejecución del ETL}

La migración se ejecutó mediante \texttt{scripts/02\_migrate\_and\_backup.sh},
que orquesta pgloader con schema de staging, precheck, carga transaccional,
postcheck, reajuste de secuencias, \texttt{ANALYZE} y limpieza. La salida real
se reproduce a continuación. Se observa el congelamiento de escrituras, el
respaldo del origen, la carga pgloader a \texttt{stg}, la validación de
precheck, la carga en una sola transacción, el reajuste de secuencias y el
respaldo cifrado del destino. El total migrado fue 1\,469 filas con 0 rechazos
y 0 errores.

\begin{lstlisting}[style=terminal,caption={Salida real de la migración con pgloader + staging},label={lst:salida-etl}]
== 1. Congelamiento de escrituras ==
Congelamiento activado en MySQL (coop_horizonte_legacy)

== 2. Respaldo del origen ==
mysqldump --single-transaction --routines --triggers --events
Archivo: backups/origin_20261120_080940.sql.gz (32 KB)
SHA-256: 6a2f3a958c4732743302299eb7e4f5d722924f749dbf6fa9764a2936428be3b3

== 3. Preparación de staging ==
psql -v ON_ERROR_STOP=1 -f sql/00_prepare_staging.sql
CREATE SCHEMA stg
ALTER DEFAULT PRIVILEGES IN SCHEMA stg GRANT ALL ON TABLES TO migration_loader

== 4. Carga pgloader a staging ==
pgloader mysql_to_stage.load
... connecting to MySQL ...
... connecting to PostgreSQL ...
types_credito:       5 rows
sucursales:          5 rows
socios:             55 rows
cuentas_ahorro:     66 rows
solicitudes_credito: 50 rows
prestamos:          32 rows
cuotas_amortizacion: 516 rows
transacciones:      740 rows
Total:            1469 rows, 0 rejects, 0 errors
ALTER SCHEMA 'coop_horizonte_legacy' RENAME TO 'stg'

== 5. Conteos origen vs staging ==
tipos_credito:       5  = 5  OK
sucursales:          5  = 5  OK
socios:             55  = 55  OK
cuentas_ahorro:     66  = 66  OK
solicitudes_credito: 50  = 50  OK
prestamos:          32  = 32  OK
cuotas_amortizacion: 516 = 516 OK
transacciones:      740 = 740 OK

== 6. Precheck ==
psql -v ON_ERROR_STOP=1 -f sql/10_precheck.sql
PRECHECK OK: nulos validados
PRECHECK OK: ENUM dentro de dominio
PRECHECK OK: sin huérfanos en claves foráneas
PRECHECK OK: rangos y fechas cero validados

== 7. Carga transaccional (20_load + 30_postcheck + 35_sequences) ==
psql --single-transaction -v ON_ERROR_STOP=1 -f sql/20_load.sql
psql --single-transaction -v ON_ERROR_STOP=1 -f sql/30_postcheck.sql
POSTCHECK OK: EXCEPT ALL bidireccional sin diferencias
psql --single-transaction -v ON_ERROR_STOP=1 -f sql/35_sequences.sql
setval OK: 8 secuencias reajustadas

== 8. Estadísticas ==
psql -U postgres -f sql/40_analyze.sql
ANALYZE OK

== 9. Limpieza de staging ==
psql -v ON_ERROR_STOP=1 -f sql/50_cleanup.sql
DROP SCHEMA stg CASCADE

== 10. Respaldo cifrado del destino ==
pg_dump -Fc coop_horizonte_core | openssl enc -aes-256-cbc -pbkdf2 -iter 200000 -salt
Archivo: backups/target_20261120_080940.dump.enc (68 KB)
SHA-256: 7017f6905da2e9b317eeb57885cb9d5ef60dff4a90c6efa6de160bf44c1c4d22

== 11. Prueba de restauración ==
Descifrado OK: 68K
pg_restore en coop_horizonte_restore
Filas de socios restauradas: 55
Restauracion: OK (coincide con el origen)

== Total migrado ==
1469 filas, 0 rechazos, 0 errores
\end{lstlisting}

\capturaAPA{fig:cap-etl}{Ejecución del ETL hacia PostgreSQL 16}%
  {captura-ejecucion-etl.png}%
  {Salida de \texttt{scripts/02\_migrate\_and\_backup.sh} con la carga de las 8 tablas y la verificación de conteos.}

La salida completa del ETL se presenta en el Anexo~C
(Figura~\ref{fig:anexo-migracion-1}).

\subsection{Respaldo cifrado del destino}

Inmediatamente después de la carga, el script generó un respaldo del destino
con \texttt{pg\_dump -Fc} y lo cifró con OpenSSL en modo AES-256-CBC, PBKDF2
con 200\,000 iteraciones y \emph{salt}. El archivo resultante fue
\texttt{backups/target\_20261120\_080940.dump.enc}, de 68 KB, con SHA-256
\texttt{7017f6905da2e9b317eeb57885cb9d5ef60dff4a90c6efa6de160bf44c1c4d22}.
La Tabla~\ref{tab:respaldos} resume los dos respaldos.

\begin{table}[H]
  \centering
  \caption{Respaldos generados durante la migración}
  \label{tab:respaldos}
  \contenidoConFuente{%
    \begin{tablaAPA}
      \begin{tabular}{@{}L{3.4cm}L{2.0cm}L{2.0cm}L{5.4cm}@{}}
        \toprule
        \textbf{Archivo} & \textbf{Tamaño} & \textbf{Formato} & \textbf{SHA-256 (prefijo)} \\
        \midrule
        \texttt{origin\_20261120\_080940.sql.gz} & 32 KB & \texttt{mysqldump} + gzip & \texttt{6a2f3a95\ldots} \\
        \texttt{target\_20261120\_080940.dump.enc} & 68 KB & \texttt{pg\_dump -Fc} + AES-256-CBC & \texttt{7017f690\ldots} \\
        \bottomrule
      \end{tabular}
    \end{tablaAPA}%
  }{Elaboración propia con datos de \texttt{evidence/02\_ejecucion\_migracion.txt}.}
\end{table}

\subsection{Prueba de restauración}

La prueba de restauración descifró el respaldo, creó la base temporal
\texttt{coop\_horizonte\_restore}, ejecutó \texttt{pg\_restore} y consultó el
conteo de socios. El resultado fue 55, idéntico al origen. La base temporal se
eliminó al finalizar. Esta prueba es fundamental: un respaldo que no se
restaura no es una garantía, es una suposición. El procedimiento se puede
repetir con cualquier respaldo archivado.

\begin{lstlisting}[style=terminal,caption={Prueba real de restauración del respaldo cifrado},label={lst:restauracion}]
== Prueba de restauracion del respaldo cifrado ==
Descifrado OK: 68K
NOTICE:  database "coop_horizonte_restore" does not exist, skipping
DROP DATABASE
CREATE DATABASE
Filas de socios restauradas: 55
Restauracion: OK (coincide con el origen)
\end{lstlisting}

\capturaAPA{fig:cap-restauracion}{Prueba de restauración del respaldo cifrado}%
  {captura-restauracion.png}%
  {Descifrado del respaldo, \texttt{pg\_restore} en la base temporal y 55 filas de socios restauradas.}

El respaldo cifrado y la restauración completa se presentan en el Anexo~C
(Figura~\ref{fig:anexo-migracion-2}).

\section{Respaldo cifrado y recuperación}

\subsection{Estrategia de respaldo}

La estrategia combina dos formatos. El origen se respalda con
\texttt{mysqldump}, que produce un archivo SQL lógico, portable y fácil de
inspeccionar; se comprime con gzip para reducir espacio. El destino se
respalda con \texttt{pg\_dump -Fc}, el formato personalizado de PostgreSQL,
que permite restaurar selectivamente y comprimir internamente. El respaldo
del destino se cifra con AES-256-CBC porque contiene datos personales y
financieros. La Tabla~\ref{tab:respaldo-controles} relaciona cada control con
su propósito.

\begin{table}[H]
  \centering
  \caption{Controles del respaldo y la recuperación}
  \label{tab:respaldo-controles}
  \contenidoConFuente{%
    \begin{tablaAPA}
      \begin{tabular}{@{}L{4.0cm}L{5.2cm}L{5.0cm}@{}}
        \toprule
        \textbf{Control} & \textbf{Implementación} & \textbf{Propósito} \\
        \midrule
        Consistencia & \texttt{mysqldump --single-transaction} & Respaldo sin bloquear la operación. \\
        Compresión & \texttt{gzip -9} en el origen; \texttt{-Fc} en el destino & Reducir tamaño y facilitar transferencia. \\
        Cifrado & AES-256-CBC, PBKDF2 200\,000 iteraciones, \emph{salt} & Proteger la confidencialidad en reposo. \\
        Integridad & SHA-256 de cada archivo & Detectar corrupción o alteración. \\
        Restauración & \texttt{pg\_restore} en base temporal y conteo & Verificar que el respaldo es utilizable. \\
        Retención & Política propuesta de 30 días diarios y 12 meses mensuales & Cumplir operación y auditoría. \\
        \bottomrule
      \end{tabular}
    \end{tablaAPA}%
  }{Elaboración propia.}
\end{table}

\subsection{Procedimiento de recuperación}

El procedimiento de recuperación tiene cinco pasos: (1) identificar el
respaldo y verificar su SHA-256; (2) descifrar con la frase de paso correcta;
(3) crear una base vacía; (4) restaurar con \texttt{pg\_restore}; y (5)
validar conteos y claves foráneas antes de abrir el servicio. En producción,
el paso 5 debe incluir una consulta de negocio (por ejemplo, saldos por
socio) y una revisión de los logs. El RTO estimado es de 30 minutos para el
volumen actual; el RPO es de 0 minutos para los datos migrados porque el
respaldo se toma en el corte.

\subsection{Gestión de claves}

En el laboratorio, la frase de paso se guarda en la variable
\texttt{BACKUP\_PASSPHRASE} del archivo \texttt{.env}, que no se versiona.
En producción, esa práctica no es suficiente: se recomienda usar un gestor de
secretos (Vault, AWS Secrets Manager, Azure Key Vault o equivalente), rotar la
frase de paso periódicamente y separar la custodia de la clave de la custodia
del respaldo. El cifrado simétrico protege el archivo, pero si la clave y el
archivo están en el mismo servidor comprometido, la protección se reduce. La
regla 3-2-1 (tres copias, dos medios, una fuera del sitio) sigue siendo una
buena referencia.

\section{Pruebas en destino}

\subsection{Resultado global}

El script \texttt{03\_verify\_all.sh} ejecutó 75 comprobaciones y obtuvo
75 \texttt{PASS} y 0 \texttt{FAIL}. La Tabla~\ref{tab:pruebas-global} resume
los grupos. Las pruebas no se limitan a verificar que el servidor responde:
comprueban el valor efectivo de los parámetros, el rechazo de conexiones sin
SSL, los privilegios de cada rol, el enmascaramiento, los conteos, las claves
foráneas, el staging, el loader y las secuencias, y el tiempo de ejecución de
las vistas.

\begin{table}[H]
  \centering
  \caption{Resumen de las pruebas en destino}
  \label{tab:pruebas-global}
  \contenidoConFuente{%
    \begin{tablaAPA}
      \begin{tabular}{@{}L{5.6cm}C{2.0cm}C{2.0cm}L{4.4cm}@{}}
        \toprule
        \textbf{Grupo} & \textbf{Pruebas} & \textbf{PASS} & \textbf{Resultado} \\
        \midrule
        Parámetros de seguridad & 5 & 5 & \texttt{ssl}, \texttt{scram}, logs. \\
        \texttt{pg\_hba} y conexiones & 4 & 4 & Sin SSL rechazado; \texttt{verify-ca} OK. \\
        Roles y privilegios & 27 & 27 & 7 roles auditados (3 operación + 1 loader + 4 admin). \\
        Enmascaramiento & 3 & 3 & CI/NIT y cuentas ocultas. \\
        Conteos & 8 & 8 & 8 tablas idénticas. \\
        Integridad referencial & 10 & 10 & 9 FK, 0 huérfanos. \\
        Staging, loader y secuencias & 18 & 18 & Schema \texttt{stg}; pgloader; \texttt{setval}; \texttt{ANALYZE}. \\
        \midrule
        \textbf{Total} & \textbf{75} & \textbf{75} & \textbf{0 FAIL}. \\
        \bottomrule
      \end{tabular}
    \end{tablaAPA}%
  }{Elaboración propia con datos de \texttt{evidence/03\_pruebas\_y\_hardening.txt}.}
\end{table}

\capturaAPA{fig:cap-pruebas}{Resumen de la verificación en destino}%
  {captura-pruebas-resumen.png}%
  {Salida final de \texttt{scripts/03\_verify\_all.sh}: \texttt{PASS: 75} y \texttt{FAIL: 0}.}

Las capturas completas de las pruebas se presentan en el Anexo~C
(Figuras~\ref{fig:anexo-pruebas-1} a~\ref{fig:anexo-pruebas-3}).

\subsection{Conteos origen--destino}

La Tabla~\ref{tab:comparativa-conteos} compara los conteos. Las 8 filas
coinciden exactamente. La coincidencia no demuestra por sí sola que los
valores sean correctos, pero junto con la verificación de claves foráneas y
la restauración del respaldo proporciona una evidencia sólida de que la carga
fue completa.

\begin{table}[H]
  \centering
  \caption{Comparativa de conteos origen vs. destino}
  \label{tab:comparativa-conteos}
  \contenidoConFuente{%
    \begin{tablaAPA}
      \begin{tabular}{@{}L{6.0cm}R{2.4cm}R{2.4cm}C{2.2cm}@{}}
        \toprule
        \textbf{Tabla} & \textbf{Origen} & \textbf{Destino} & \textbf{Estado} \\
        \midrule
        \texttt{sucursales} & 5 & 5 & OK \\
        \texttt{socios} & 55 & 55 & OK \\
        \texttt{cuentas\_ahorro} & 66 & 66 & OK \\
        \texttt{tipos\_credito} & 5 & 5 & OK \\
        \texttt{solicitudes\_credito} & 50 & 50 & OK \\
        \texttt{prestamos} & 32 & 32 & OK \\
        \texttt{cuotas\_amortizacion} & 516 & 516 & OK \\
        \texttt{transacciones} & 740 & 740 & OK \\
        \midrule
        \textbf{Total} & \textbf{1\,469} & \textbf{1\,469} & \textbf{OK} \\
        \bottomrule
      \end{tabular}
    \end{tablaAPA}%
  }{Elaboración propia con datos de \texttt{evidence/03\_pruebas\_y\_hardening.txt}.}
\end{table}

\subsection{Integridad referencial}

Se verificaron las 9 claves foráneas con consultas de huérfanos. La
Tabla~\ref{tab:fks} muestra las relaciones y el resultado. Todas devolvieron
0. La consulta genérica es \texttt{SELECT COUNT(*) FROM hija h LEFT JOIN padre
p ON h.fk = p.pk WHERE h.fk IS NOT NULL AND p.pk IS NULL}. Si el resultado es
distinto de cero, la clave foránea no se puede declarar o la carga tiene un
error.

\begin{table}[H]
  \centering
  \caption{Claves foráneas verificadas}
  \label{tab:fks}
  \contenidoConFuente{%
    \begin{tablaAPA}
      \begin{tabular}{@{}L{4.4cm}L{4.4cm}R{2.2cm}@{}}
        \toprule
        \textbf{Tabla hija} & \textbf{Tabla padre} & \textbf{Huérfanos} \\
        \midrule
        \texttt{cuentas\_ahorro.socio\_id} & \texttt{socios.id} & 0 \\
        \texttt{solicitudes\_credito.socio\_id} & \texttt{socios.id} & 0 \\
        \texttt{solicitudes\_credito.tipo\_credito\_id} & \texttt{tipos\_credito.id} & 0 \\
        \texttt{prestamos.solicitud\_id} & \texttt{solicitudes\_credito.id} & 0 \\
        \texttt{prestamos.socio\_id} & \texttt{socios.id} & 0 \\
        \texttt{cuotas\_amortizacion.prestamo\_id} & \texttt{prestamos.id} & 0 \\
        \texttt{transacciones.socio\_id} & \texttt{socios.id} & 0 \\
        \texttt{transacciones.cuenta\_ahorro\_id} & \texttt{cuentas\_ahorro.id} & 0 \\
        \texttt{transacciones.prestamo\_id} & \texttt{prestamos.id} & 0 \\
        \bottomrule
      \end{tabular}
    \end{tablaAPA}%
  }{Elaboración propia con datos de \texttt{evidence/03\_pruebas\_y\_hardening.txt}.}
\end{table}

\subsection{Rendimiento de consultas críticas}

Se midieron las dos vistas principales. La Tabla~\ref{tab:rendimiento}
distingue el tiempo de pared (\emph{wall}) del tiempo de servidor. El primero
incluye el costo de \texttt{docker compose exec}, la autenticación, la
transferencia y el formateo; el segundo es el \texttt{Execution Time} de
\texttt{EXPLAIN ANALYZE} y refleja el trabajo real del motor. Para
\texttt{vw\_cartera\_mora}, el servidor tardó 0.236 ms; para
\texttt{vw\_resumen\_saldos\_socio}, 0.317 ms.

\begin{table}[H]
  \centering
  \caption{Rendimiento de las vistas críticas}
  \label{tab:rendimiento}
  \contenidoConFuente{%
    \begin{tablaAPA}
      \begin{tabular}{@{}L{5.0cm}R{2.6cm}R{2.6cm}R{2.8cm}@{}}
        \toprule
        \textbf{Consulta} & \textbf{MySQL wall (ms)} & \textbf{PG wall (ms)} & \textbf{PG servidor (ms)} \\
        \midrule
        \texttt{vw\_cartera\_mora} & 124 & 143 & 0.236 \\
        \texttt{vw\_resumen\_saldos\_socio} & 108 & 137 & 0.317 \\
        \bottomrule
      \end{tabular}
    \end{tablaAPA}%
  }{Elaboración propia con datos de \texttt{evidence/03\_pruebas\_y\_hardening.txt}.}
\end{table}

El plan de \texttt{vw\_cartera\_mora} usa un \texttt{Bitmap Index Scan} sobre
\texttt{idx\_prestamo\_estado} para localizar los préstamos morosos, luego un
\texttt{Hash Join} con socios y cuotas, y un \texttt{GroupAggregate}. El plan
de \texttt{vw\_resumen\_saldos\_socio} usa \texttt{Merge Left Join} e
\texttt{Incremental Sort}. Para el volumen actual, los \texttt{Seq Scan} sobre
tablas pequeñas son razonables; a mayor escala, el optimizador podría cambiar
a índices si las estadísticas y la selectividad lo justifican. La lección es
que el rendimiento debe medirse con \texttt{EXPLAIN (ANALYZE, BUFFERS)} y no
suponerse por el tamaño de la tabla.

\section{Confidencialidad y escalabilidad}

\subsection{Clasificación de la información}

La Tabla~\ref{tab:clasificacion} clasifica los datos tratados. Los datos
personales y financieros son confidenciales o restringidos. La clasificación
guía las medidas: los datos restringidos solo se exponen a través de vistas
enmascaradas; los confidenciales se cifran en tránsito y en reposo; los
internos se protegen con roles; y los públicos pueden compartirse sin
restricción adicional.

\begin{table}[H]
  \centering
  \caption{Clasificación de la información}
  \label{tab:clasificacion}
  \contenidoConFuente{%
    \begin{tablaAPA}
      \begin{tabular}{@{}L{2.8cm}L{5.0cm}L{6.0cm}@{}}
        \toprule
        \textbf{Nivel} & \textbf{Ejemplos} & \textbf{Control principal} \\
        \midrule
        Restringido & CI/NIT, número de cuenta, teléfono, correo, dirección. & Enmascaramiento en \texttt{vw\_socios\_audit}; sin acceso directo para \texttt{auditor\_ro}. \\
        Confidencial & Saldos, montos de crédito, cuotas, transacciones. & TLS, roles, respaldo cifrado, auditoría. \\
        Interno & Catálogos de sucursales y tipos de crédito. & Roles y registros de acceso. \\
        Público & Documentación, scripts, informe. & Sin restricción adicional. \\
        \bottomrule
      \end{tabular}
    \end{tablaAPA}%
  }{Elaboración propia.}
\end{table}

\subsection{Controles de confidencialidad}

Los controles se organizan en cuatro momentos: en tránsito, en reposo, en uso
y en auditoría. En tránsito, TLS 1.2+ con \texttt{hostssl} y
\texttt{scram-sha-256}. En reposo, AES-256-CBC para respaldos y permisos
\texttt{600} para la clave privada. En uso, roles de mínimo privilegio y
vistas enmascaradas. En auditoría, registros de conexión, desconexión y DDL,
más la separación del rol \texttt{auditor\_ro}. La combinación es más fuerte
que cualquier control aislado: si un atacante obtiene una credencial de
aplicación, no puede leer la tabla de socios; si obtiene un respaldo, no puede
descifrarlo sin la frase de paso; si intenta conectarse sin TLS, el servidor
lo rechaza.

\subsection{Escalabilidad}

El diseño actual es adecuado para el volumen de la cooperativa, pero la
escalabilidad requiere un plan. La Tabla~\ref{tab:escalabilidad} propone
medidas por horizonte. A corto plazo, se recomienda añadir un \emph{pool} de
conexiones (PgBouncer), monitorear consultas lentas y ajustar
\texttt{autovacuum}. A mediano plazo, particionar \texttt{transacciones} por
fecha y crear índices parciales para los estados más consultados. A largo
plazo, evaluar réplicas de lectura para reportes y una arquitectura de alta
disponibilidad con \emph{failover}.

\begin{table}[H]
  \centering
  \caption{Medidas de escalabilidad por horizonte}
  \label{tab:escalabilidad}
  \contenidoConFuente{%
    \begin{tablaAPA}
      \begin{tabular}{@{}L{2.4cm}L{6.2cm}L{5.6cm}@{}}
        \toprule
        \textbf{Horizonte} & \textbf{Medida} & \textbf{Beneficio} \\
        \midrule
        Corto plazo & PgBouncer, monitoreo de consultas lentas, ajuste de \texttt{autovacuum}. & Menos conexiones, mejor uso de recursos. \\
        Mediano plazo & Particionamiento de \texttt{transacciones} por rango de fechas; índices parciales. & Consultas históricas más rápidas; mantenimiento por partición. \\
        Largo plazo & Réplicas de lectura para reportes; alta disponibilidad con \emph{failover}. & Escalabilidad horizontal y continuidad. \\
        Transversal & Capacitación, documentación y pruebas de carga periódicas. & Operación sostenible y decisiones basadas en datos. \\
        \bottomrule
      \end{tabular}
    \end{tablaAPA}%
  }{Elaboración propia.}
\end{table}

La escalabilidad no es solo agregar hardware: es evitar cuellos de botella
conocidos. Un \emph{pool} mal configurado puede ser peor que no tenerlo; un
índice innecesario ralentiza las escrituras; una réplica retrasada puede
servir datos obsoletos. Por eso las medidas se proponen por horizonte y con
métricas de validación, no como una lista de deseos.
